\subsection{Hillenbrand's explanation: long-run Fed guidance}

Hillenbrand's own interpretation is that FOMC announcements are the moments at which the Fed
reveals information about \emph{where rates are going in the long run} --- what he calls
\textbf{long-run Fed guidance}. The premise is an information asymmetry: by continuously
observing how the economy responds to the rates it sets, the Fed accumulates an estimate of
the natural rate that markets cannot form as precisely on their own. Investors learn that
estimate at announcements and reprice long-duration bonds accordingly. Because the object
being revised is an \emph{endpoint} rather than a cyclical position, the revision is
permanent, and permanent revisions cumulate rather than cancel --- which is exactly the
pattern in Figure~\ref{fig:fact}.

He describes the guidance as operating through two channels:

\begin{itemize}\setlength\itemsep{0.15em}
\item \textbf{Implicit, before 2012.} When the Fed moves the funds rate, the move itself
reveals its reading of where the natural rate sits. Consistent with this, the meetings at
which the target changed --- roughly 45\% of meetings in his pre-2012 sample --- account for
almost the entire decline in long yields over that period.
\item \textbf{Explicit, from 2012.} The SEP dot plot publishes each participant's longer-run
federal funds projection, communicating the Fed's endpoint view directly.
\end{itemize}

\noindent Four pieces of evidence are offered that the object being moved is a long-run
\emph{real} endpoint rather than the cycle, inflation, or a risk premium.

\begin{enumerate}\setlength\itemsep{0.2em}
\item \textbf{Far-forward rates.} The pattern holds for the 5-year, 5-year-forward rate, where
the current policy stance has largely washed out. A cyclical explanation should not survive
that far out the curve.
\item \textbf{Real rather than inflation.} Since 1999 the decline is essentially all in TIPS
yields, with breakeven inflation roughly flat --- so it is $\Delta r^{*}$, not $\Delta\pi^{*}$.
\item \textbf{The dot plot links the two legs directly.} Regressing window changes in
long-term real yields on revisions to the Fed's longer-run projection, a 100 bp cut in the
Fed's own long-run forecast is associated with more than a 70 bp fall in long-term real
yields --- a pass-through large enough to be about the endpoint rather than the near path.
\item \textbf{Persistence.} The series built only from window changes tracks the low-frequency
movement of the ten-year closely, while the non-window series is transitory and offsets over
time. Endpoint news should not reverse; cyclical or liquidity-driven news should.
\end{enumerate}

\noindent In terms of the identity developed in \S\ref{sec:identity}, the claim is a strong
one: in FOMC windows $\Delta y^{(n)} \approx \Delta r^{*}$, with the inflation anchor, the
policy-stance terms and the term premium all contributing approximately nothing. Each of
those three exclusions is a separate assumption, and the rest of this report is largely about
which of them survive.
